\documentclass[11pt]{article}
\usepackage[margin=1in]{geometry}
\usepackage[T1]{fontenc}
\usepackage[utf8]{inputenc}
\usepackage{hyperref}
\usepackage{enumitem}
\setlist[itemize]{leftmargin=1.5em}
\setlist[enumerate]{leftmargin=1.7em}

\title{Literature Review: Space Debris Collision Avoidance}
\author{}
\date{}

\begin{document}
\maketitle

\section*{Team Details}
\begin{itemize}
\item Mayank Goyal (Team Leader), Roll No: 24293916114, Email: mayankgoyal18052005@ce.du.ac.in
\item Abeer bartaria, Roll No: 23293916061, Email: abeerbartaria@ce.du.ac.in
\item Ansh Kumar, Roll No: 24293916052, Email: Anshk3671@gmail.com
\item Vansh Saharawat, Roll No: 24293916089, Email: vanshsahara@ce.du.ac.in
\end{itemize}

\section*{Scope}
This review covers conjunction assessment and collision avoidance in Earth orbit, aligned with our EDA domain.

\section*{Literature Collection (5 Works)}

\subsection*{1) Chan (2000)}
Source: Journal of Guidance, Control, and Dynamics\\
DOI: \url{https://doi.org/10.2514/2.4611}
\begin{itemize}
\item Problem statement: How to compute collision risk rigorously for close approaches.
\item Dataset used: Cataloged space objects and encounter uncertainty states.
\item Methods applied: Analytical probability of collision using encounter geometry and covariance.
\item Key findings: Probability-based thresholds are better than simple geometric screening rules.
\end{itemize}

\subsection*{2) Newman et al. (2009)}
Source: Acta Astronautica\\
DOI: \url{https://doi.org/10.1016/j.actaastro.2009.10.005}
\begin{itemize}
\item Problem statement: How to run robotic mission conjunction assessment in operations.
\item Dataset used: U.S. tracking catalog conjunction products plus mission trajectory data.
\item Methods applied: Multi-stage screening, escalation, analyst review, and maneuver planning.
\item Key findings: Layered pipelines reduce false alerts and support practical mission decisions.
\end{itemize}

\subsection*{3) Flohrer et al. (2009)}
Source: Advances in Space Research\\
DOI: \url{https://doi.org/10.1016/j.asr.2009.04.012}
\begin{itemize}
\item Problem statement: How to identify and assess high-risk conjunction events for ESA missions.
\item Dataset used: Daily screenings of tracked objects with orbit prediction and uncertainty data.
\item Methods applied: Automated filtering, manual refinement, trend checks, and decision gates.
\item Key findings: Staged assessment improves safety while limiting unnecessary maneuvers.
\end{itemize}

\subsection*{4) Braun et al. (2016)}
Source: CEAS Space Journal\\
DOI: \url{https://doi.org/10.1007/s12567-016-0119-3}
\begin{itemize}
\item Problem statement: How to scale operational collision support across many missions.
\item Dataset used: Operational conjunction streams, OD updates, and mission constraints.
\item Methods applied: Standardized support workflow and mission-tailored risk handling.
\item Key findings: Institutional support improves consistency, but upstream uncertainty remains a major limiter.
\end{itemize}

\subsection*{5) Weigel et al. (2024)}
Source: Journal of Guidance, Control, and Dynamics\\
DOI: \url{https://doi.org/10.2514/1.G008245}
\begin{itemize}
\item Problem statement: How to make covariance models realistic enough for robust CA decisions.
\item Dataset used: Historical mission orbit-estimation and covariance behavior data.
\item Methods applied: Consider-covariance calibration and comparative risk-ranking experiments.
\item Key findings: Better covariance realism improves event ranking and maneuver timing quality.
\end{itemize}

\section*{Comparative Analysis with Our EDA}

\subsection*{Our EDA snapshot}
\begin{itemize}
\item SATCAT cleaned: 68,330 objects.
\item UCS cleaned: 7,563 satellites.
\item Kelvins labels: 100 records, with long trajectory subsets (deb\_train, deb\_test, sat).
\item UCS analysis shows LEO-heavy distribution and Communications-dominant mission purpose.
\item Launch trend rises sharply in recent years, especially around 2020 to 2022.
\end{itemize}

\subsection*{Similarities with literature}
\begin{itemize}
\item Literature and our EDA both indicate strong growth in orbital traffic and conjunction burden.
\item LEO concentration in our UCS data matches operational studies where LEO dominates conjunction workload.
\item Data quality and uncertainty are central in both; our cleaning steps already show mixed formats and missingness that impact downstream risk scoring.
\end{itemize}

\subsection*{Differences from literature}
\begin{itemize}
\item Literature mostly uses mission-grade conjunction products (for example CDMs and refined covariance chains), while our current datasets are catalog-level and benchmark-oriented.
\item Literature includes maneuver authority and mission constraints; our current pipeline is pre-decision and analytics-focused.
\item Recent covariance-optimization methods require historical orbit-determination streams not present in our current data bundle.
\end{itemize}

\section*{Research Gaps and Future Work}

\subsection*{Gap 1: Public end-to-end CA benchmarks are limited}
Most strong systems are agency internal. Reproducible public benchmarks that cover screening, ranking, and maneuver recommendation are still limited.

\subsection*{Gap 2: Screening and maneuver planning are often disconnected}
Many workflows optimize each stage separately. A unified objective from first filter to final maneuver can reduce inconsistency.

\subsection*{Gap 3: Uncertainty treatment is still uneven}
Catalog-level studies often under-model covariance realism. Operational papers show uncertainty quality strongly changes decisions.

\subsection*{Proposed next steps for our project}
\begin{enumerate}
\item Build a two-stage CA pipeline on top of our processed SATCAT and UCS artifacts: fast screening plus refined probability scoring.
\item Add uncertainty-aware modeling inspired by modern covariance calibration work.
\item Evaluate with operationally meaningful metrics: false alert rate, unnecessary maneuver rate, and missed-danger events.
\item Publish a reproducible benchmark protocol with fixed scenario splits and reporting templates.
\end{enumerate}

\section*{References}
\begin{enumerate}
\item Chan, F. K. (2000). Probability of Collision Between Space Objects. Journal of Guidance, Control, and Dynamics. \url{https://doi.org/10.2514/2.4611}
\item Newman, L. K., et al. (2009). The NASA robotic conjunction assessment process: Overview and operational experiences. Acta Astronautica. \url{https://doi.org/10.1016/j.actaastro.2009.10.005}
\item Flohrer, T., Krag, H., and Klinkrad, H. (2009). ESA's process for the identification and assessment of high-risk conjunction events. Advances in Space Research. \url{https://doi.org/10.1016/j.asr.2009.04.012}
\item Braun, V., et al. (2016). Operational support to collision avoidance activities by ESA's space debris office. CEAS Space Journal. \url{https://doi.org/10.1007/s12567-016-0119-3}
\item Weigel, M., et al. (2024). Optimizing Consider-Covariance Models for Realistic Orbit Errors and Application to Collision Avoidance. Journal of Guidance, Control, and Dynamics. \url{https://doi.org/10.2514/1.G008245}
\end{enumerate}

\end{document}
